\documentclass[11pt,a4paper]{article}
\usepackage[margin=2.15cm]{geometry}
\usepackage[T1]{fontenc}
\usepackage{lmodern}
\usepackage{amsmath,amssymb,mathtools}
\usepackage{enumitem}
\usepackage{fancyhdr}
\usepackage{booktabs}
\usepackage{tikz}
\usepackage[hidelinks]{hyperref}
\usepackage{microtype}

\setlength{\parindent}{0pt}
\setlength{\parskip}{0.45em}
\setlist[enumerate]{leftmargin=*,itemsep=0.4em,topsep=0.3em}
\pagestyle{fancy}
\setlength{\headheight}{14pt}
\fancyhf{}
\lhead{Physics Practice Examination}
\rhead{\thepage}
\renewcommand{\headrulewidth}{0.4pt}

\newcommand{\problem}[3]{%
  \subsection*{Problem #1 \;\textnormal{\small[#2]}}\label{prob:#3}%
}
\newcommand{\unit}[1]{\,\mathrm{#1}}

\newenvironment{parts}{\begin{enumerate}[label=\textbf{\alph*)}]}{\end{enumerate}}

\title{\vspace{-1.4cm}\bfseries Physics Practice Examination\\[2pt]
\large Finnish Upper Secondary Syllabus (LOPS 2021), Modules FY1--FY8}
\author{Ten Problems, Weighted Towards FY5--FY8}
\date{}

\begin{document}
\maketitle
\vspace{-1.6em}

\begin{center}
\fbox{\parbox{0.92\textwidth}{
\textbf{Suggested time: 4 hours.} \quad 10 problems, 6 points each, 60 points
total. Calculators and the official formula booklet (\emph{MAOL}) are
permitted. Full solutions with units are expected; numerical answers to
2--3 significant figures unless the problem asks for an exact value.
Problems are \emph{not} in order of difficulty.\\[3pt]
\textbf{Course mapping:} P1 (FY1) $\cdot$ P2 (FY2) $\cdot$ P3 (FY3) $\cdot$
P4 (FY4) $\cdot$ P5, P6 (FY5) $\cdot$ P7 (FY6) $\cdot$ P8, P9 (FY7) $\cdot$
P10 (FY8). Problems 5--10 carry the weight of the higher courses.
}}
\end{center}

%======================================================================
\problem{1}{FY1 --- Measurement, graphical models, uncertainty}{fy1}

A student determines the acceleration of free fall $g$ with a simple pendulum.
The bob has radius $r$ and hangs on a thread of negligible mass; the
\emph{effective} length of the pendulum, the quantity that enters the formula
for the period, is thread length plus bob radius. Little time is available, so
she makes \emph{one} run at each of five thread lengths and obtains:

\begin{center}
\begin{tabular}{cc}
\toprule
thread length $l\unit{(m)}$ & period $T\unit{(s)}$\\
\midrule
0.200 & 0.914\\
0.400 & 1.287\\
0.600 & 1.573\\
0.800 & 1.813\\
1.000 & 2.025\\
\bottomrule
\end{tabular}
\end{center}

The periods were measured with a photogate that starts and stops the timer
itself, so the random timing error is at the millisecond level. The angular
amplitude was $20.0^\circ$ in every run, with a reading uncertainty of
$\pm0.5^\circ$.

\begin{parts}
\item Explain why plotting $T^2$ against $l$, rather than $T$ against $l$,
      gives a straight line, and state what theory predicts for the slope and
      the intercept of that line (keep $r$ as a symbol).
\item Determine $g$ and its uncertainty from the scatter of the points about
      the fitted line. Quote $g\pm\Delta g$ with a sensible number of
      significant figures.
\item Determine the bob radius $r$ from the intercept, and judge whether the
      intercept is significantly different from zero.
\item Your value turns out to be about $1.4\%$ below $9.81\unit{m/s^2}$. Show
      quantitatively that the finite angular amplitude accounts for this, and
      compute the corrected result. Then state, in the language of
      measurement error, what kind of error this is, what it did to the
      estimate of $g$, and whether it is visible in the uncertainty of
      part (b). Finally, say what the experimenter should have done instead,
      given that the photogate is already available.
\end{parts}

%======================================================================
\problem{2}{FY2 --- Energy production, power, efficiency}{fy2}

A wind farm is planned with turbines of rotor diameter $D=120\unit{m}$ in air
of density $\rho=1.20\unit{kg/m^3}$. The wind is described by a Rayleigh
distribution, i.e.\ it is characterised by its mean speed $\bar v$, here
$\bar v=7.5\unit{m/s}$. In such a distribution the mean of the \emph{cube} of
the speed exceeds the cube of the mean speed:
$\overline{v^3}=\tfrac{6}{\pi}\,\bar v^{\,3}$. Each turbine delivers a rated
electrical power $P_{\text{rated}}=2.50\unit{MW}$ above a certain wind speed,
and the farm's annual capacity factor is $0.35$. Finnish electricity
consumption is about $90\unit{TWh}$ per year, and coal condensing power emits
about $0.90\unit{kg}$ of CO$_2$ per kWh.

\begin{parts}
\item Derive the power flowing through the rotor disc,
      $P_w=\tfrac12\rho A v^3$, from dimensional and physical reasoning, and
      evaluate it at $v=12\unit{m/s}$. Why can no turbine extract all of it?
      State the Betz limit and compare it with the fraction your turbine
      actually achieves at $12\unit{m/s}$.
\item A turbine's average output is
      $\overline{P}=\tfrac12\rho A\,\overline{v^3}\,C_p$, where $C_p$ is the
      power coefficient. Compute first $\tfrac12\rho A\bar v^{\,3}$, then the
      correct average available power, and hence $C_p$. Check that $C_p$ is
      below the Betz limit, and explain in one sentence why using
      $\bar v^{\,3}$ instead of $\overline{v^3}$ here would be an error, and
      in which direction.
\item Compute the annual electrical energy of one turbine in GWh, the number
      of turbines needed for $1.0\%$ of Finnish consumption, and the CO$_2$
      avoided per turbine-year relative to coal power.
\item On one site the wind alternates between $5\unit{m/s}$ and
      $10\unit{m/s}$, half the time each, so the mean speed is
      $7.5\unit{m/s}$. Show that this site yields more average power than a
      site with a perfectly steady $7.5\unit{m/s}$ wind, and give the ratio.
      Then explain why a real turbine nevertheless does not deliver its
      maximum power at the highest wind speeds.
\end{parts}

%======================================================================
\problem{3}{FY3 --- Thermodynamics, ideal gas, heat}{fy3}

A hot-air balloon has envelope volume $V=2.2\cdot10^{3}\unit{m^3}$. The air
outside is at $T_0=20.0^\circ$C and pressure $p_0=101.3\unit{kPa}$; the molar
mass of air is $M=29.0\unit{g/mol}$. The envelope material has mass
$180\unit{kg}$; basket and burner together have mass $90\unit{kg}$. The
pressure inside equals the pressure outside at the open mouth, so the internal
density follows from the ideal gas law. Take $c_p=1.005\unit{kJ/(kg\,K)}$ for
air and $R=8.314\unit{J/(mol\,K)}$.

\begin{parts}
\item Compute the density of the outside air.
\item What must the average temperature of the air inside the balloon be for
      the balloon to carry a payload of $350\unit{kg}$?
\item The fabric is rated for at most $120^\circ$C. What is the largest
      payload the balloon can lift?
\item In steady flight the burner heats air entering through the mouth at
      $0.45\unit{m^3/s}$ (measured at the internal temperature) and expels it
      at the top. Estimate the burner power needed to hold the temperature
      found in part (b). State the main approximation, and say qualitatively
      why the true figure is somewhat larger.
\end{parts}

%======================================================================
\problem{4}{FY4 --- Forces, Newton's laws, energy, momentum}{fy4}

Two blocks are connected by a light, inextensible string over a frictionless
ideal pulley. Block A, $m_A=4.00\unit{kg}$, slides on a plane inclined at
$\theta=30.0^\circ$; block B, $m_B=3.00\unit{kg}$, hangs vertically. The
coefficients of friction between A and the plane are $\mu_s=0.25$ and
$\mu_k=0.20$. The system is released from rest, and B falls a distance
$h=0.80\unit{m}$ before hitting the floor.

\begin{parts}
\item Show that the system really does start to move as described, with block
      B going down, and compute the acceleration of the blocks and the
      tension in the string while B falls.
\item Compute the speed with which B strikes the floor, and verify the answer
      with an energy balance for the whole system, listing every term.
\item After B has landed the string goes slack. Determine how much further
      block A slides up the plane before stopping, and whether A then stays
      put. Justify your answer with the numbers.
\item A now slides back down. Show that it returns to the position it had at
      the instant B landed, and find its speed there. At that instant the
      string snaps taut, jerking B (still lying on the floor) into motion.
      Determine the common speed of the blocks immediately afterwards and the
      impulse transmitted by the string, and show that the fraction of
      kinetic energy lost is $m_B/(m_A+m_B)$. Then determine what happens in
      the next instant, and explain why the system does not simply continue
      to move together.
\end{parts}

%======================================================================
\problem{5}{FY5 --- Circular motion, gravitation, oscillations}{fy5}

The International Space Station orbits at an altitude of $400\unit{km}$. Use
$M_E=5.97\cdot10^{24}\unit{kg}$, $R_E=6.371\cdot10^{6}\unit{m}$ and
$G=6.674\cdot10^{-11}\unit{N\,m^2/kg^2}$. Ignore air drag except where asked.

\begin{parts}
\item Compute the orbital speed, the orbital period, and the local
      gravitational field strength at the station's altitude.
\item Show that for a circular orbit of radius $r$ the kinetic energy of a
      body of mass $m$ equals $-\tfrac12U$, where $U$ is its gravitational
      potential energy with $U\to0$ at infinity. The station has mass
      $4.2\cdot10^{5}\unit{kg}$; compute its total mechanical energy. Then
      compute the escape speed from $400\unit{km}$ and explain, using
      $K=-\tfrac12U$, why it exceeds the orbital speed by exactly
      $\sqrt2$.
\item How much work must be done \emph{on} the station to lower it from
      $400\unit{km}$ to $399\unit{km}$? Hence estimate the changes in $v$ and
      in the period, and explain the sign of the change in $v$ --- in
      particular why a decelerating manoeuvre makes the station go
      \emph{faster}.
\item A crew member in a planned rotating habitat of angular speed
      $\omega=0.256\unit{rad/s}$ sleeps with head towards the rotation axis,
      feet against the outer wall at radius $R_s=50.0\unit{m}$. Compute the
      artificial gravity at the feet, and show that it is $3.6\%$ weaker at
      the head of a person $1.80\unit{m}$ tall. Why is this gradient a
      problem, and what does it imply about the size of a healthy rotating
      habitat?
\item A $3.00\unit{m}$ pendulum hangs in that habitat. What is its
      small-oscillation period, and how does it compare with the same
      pendulum on Earth?
\end{parts}

%======================================================================
\problem{6}{FY5 --- Sound, resonance, superposition}{fy5s}

A tuning fork marked $440\unit{Hz}$ produces sound whose displacement
amplitude at a nearby point is $0.500\unit{mm}$; the speed of sound in air is
$343\unit{m/s}$. In a resonance-tube experiment with a closed pipe (one open
end) the first two resonances of this fork occur at effective pipe lengths
$l_1=0.182\unit{m}$ and $l_2=0.572\unit{m}$.

\begin{parts}
\item Compute the maximum speed and maximum acceleration of the air
      molecules, and the wavelength of the sound. Show that
      $\lambda=2(l_2-l_1)$, and compute the speed of sound from the
      measurement.
\item The first resonance is found at a tube length \emph{smaller} than the
      $\lambda/4$ one would predict for an ideal pipe. Explain the physical
      reason and determine the end correction $\Delta l$. Using
      $\Delta l\approx0.6\,r$ for a flanged open end, estimate the tube
      radius. Why does $\Delta l$ cancel out of $\lambda=2(l_2-l_1)$ but not
      out of $\Delta l=l_{\text{ideal}}-l_1$?
\item The same fork is sounded together with a second fork and beats of
      $4.0\unit{Hz}$ are heard. What are the possible frequencies of the
      second fork? Describe an experiment that decides between them, and
      state how many beats are heard in $5.0\unit{s}$ for
      $f_2=444\unit{Hz}$.
\item The experiment is repeated in a room filled with carbon dioxide and the
      first resonance is then at $l_1=0.140\unit{m}$. Assuming the same end
      correction, determine the speed of sound in CO$_2$ and compare it with
      the value predicted from $v=\sqrt{\gamma RT/M}$ using
      $\gamma_{\text{CO}_2}=1.30$, $M_{\text{CO}_2}=44.0\unit{g/mol}$ and the
      corresponding air values $\gamma=1.40$, $M=29.0\unit{g/mol}$.
      Would the beat frequency of part (c) change in CO$_2$?
\end{parts}

%======================================================================
\problem{7}{FY6 --- DC circuits and capacitors}{fy6}

In the circuit below, $E=12.0\unit{V}$, the battery's internal resistance is
$r=3.0\unit{\Omega}$, $R_1=18\unit{\Omega}$, $R_2=22\unit{\Omega}$,
$R_3=15\unit{\Omega}$ and $C=2.2\unit{\mu}$F. The capacitor has negligible
leakage, and the circuit has been connected for a long time.

\begin{center}
\begin{tikzpicture}[scale=0.95,>=latex,every node/.style={font=\small}]
  \draw (0,0) -- (0,1.75);
  \draw[line width=1.2pt] (-0.32,2.10) -- (0.32,2.10);
  \draw[line width=2.4pt] (-0.16,1.75) -- (0.16,1.75);
  \draw (0,2.10) -- (0,4.20);
  \node[left] at (-0.40,1.92) {$E,r$};
  \draw (0,4.20) -- (4,4.20) -- (4,3.65);
  \node[draw,minimum width=1.2cm,minimum height=0.45cm] at (4,3.42) {$R_3$};
  \draw (4,3.20) -- (4,2.70);
  \fill (4,2.70) circle (2pt);
  \draw (4,2.70) -- (9,2.70);
  \node[draw,minimum width=1.2cm,minimum height=0.45cm] at (4,1.35) {$R_1$};
  \draw (4,2.70) -- (4,1.58); \draw (4,1.12) -- (4,0);
  \node[draw,minimum width=1.2cm,minimum height=0.45cm] at (6.5,1.35) {$R_2$};
  \draw (6.5,2.70) -- (6.5,1.58); \draw (6.5,1.12) -- (6.5,0);
  \node[draw,minimum width=0.9cm,minimum height=0.45cm] at (9,1.35) {$C$};
  \draw (9,2.70) -- (9,1.58); \draw (9,1.12) -- (9,0);
  \draw (0,0) -- (9,0);
\end{tikzpicture}\\[1pt]
{\small The pair $R_1\parallel R_2$ is in series with $R_3$ and the internal
resistance; the capacitor is in parallel with that pair.}
\end{center}

\begin{parts}
\item Determine the total current, the terminal voltage of the battery, the
      current in each of $R_1$ and $R_2$, and the charge and energy stored in
      the capacitor.
\item Compute the power dissipated in every resistor, verify energy
      conservation, and determine the circuit's efficiency, defined as the
      fraction of the chemical power that is dissipated in the external
      resistors.
\item A voltmeter of finite resistance $R_V$ is connected across the parallel
      pair. Derive how the reading depends on $R_V$, evaluate it for
      $R_V=200\unit{\Omega}$ and $R_V=30\unit{\Omega}$, and explain why these
      readings are systematically too small. What would you require of $R_V$
      for a $1\%$ error?
\item An LED with forward voltage $2.0\unit{V}$ and rated current
      $20\unit{mA}$ is to be run from a separate $9.0\unit{V}$ battery of
      negligible internal resistance. Explain why a diode's
      current--voltage characteristic makes a series resistor necessary, and
      determine its resistance and the power it dissipates. What happens if
      the LED is connected the wrong way round, and why does the current fall
      as such a battery discharges?
\end{parts}

%======================================================================
\problem{8}{FY7 --- Induction, generators, transformers}{fy7}

\begin{parts}
\item A $100$-turn coil of cross-sectional area $1.22\cdot10^{-2}\unit{m^2}$
      rotates at $3000$ revolutions per minute in a uniform field of
      $0.85\unit{T}$. Derive the induced emf as a function of time, and
      compute the peak and rms values. Explain why the peak emf exceeds the
      average emf over a quarter turn by exactly the factor you find, and
      state the physical reason for the factor.
\item That generator feeds a $50\unit{MW}$ load. Electricity is transmitted
      through a line of total resistance $10\unit{\Omega}$. Compute the line
      loss both without a transformer (applying the coil's own rms voltage)
      and with the voltage raised to $110\unit{kV}$, and give the lost
      fraction of the power in each case. Why can a transformer not be used
      with a DC source? Why is the transmission voltage a question of
      \emph{cost}, not of the transformer's efficiency, and what fixes the
      practical upper limit on it?
\end{parts}

%======================================================================
\problem{9}{FY7 --- Induction, wave optics, optical fibre}{fy7b}

\begin{parts}
\item A straight metal rod of mass $m=0.20\unit{kg}$ and length
      $L=0.30\unit{m}$ slides on two parallel, frictionless, perfectly
      conducting rails inclined at $30.0^\circ$ to the horizontal. The rails
      are joined at the top by a resistor $R=0.30\unit{\Omega}$, so the rod
      and resistor form a closed loop. A uniform magnetic field of magnitude
      $B=1.2\unit{T}$ is applied perpendicular to the plane of the rails. The
      rod is released from rest.
      \begin{enumerate}[label=\roman*.]
      \item Derive the terminal speed and compute it.
      \item At the terminal speed, compute the current and verify that the
            electrical power equals the rate of loss of mechanical energy.
      \item From the terminal state, an external force pushes the rod
            \emph{up} the incline at a constant $2.0\unit{m/s}$. Find the
            required force and the current, and account for the power balance
            among the three energy channels.
      \item Now suppose the field is \emph{vertical} (magnitude $1.2\unit{T}$)
            while the rails still lie in the inclined plane. Explain why a
            current \emph{does} flow (compute $\mathcal{E}=BLv\cos\theta$),
            and show that the terminal speed is larger by the factor
            $1/\cos^2\theta$ than in the perpendicular case. Why does the
            magnetic force now need a component along the rails of
            $mg\sin\theta$, and where does the extra factor of $1/\cos\theta$
            come from?
      \end{enumerate}
\item Light of vacuum wavelength $\lambda=500\unit{nm}$ enters a
      $1.00\unit{m}$ long silica fibre whose core has index $n_c=1.46$ and
      cladding $n_{cl}=1.45$.
      \begin{enumerate}[label=\roman*.]
      \item Determine the critical angle at the core--cladding boundary, and
            the greatest angle to the fibre \emph{axis} at which light is
            still guided inside the core. Hence compute the numerical
            aperture $\mathrm{NA}=\sqrt{n_c^2-n_{cl}^2}$ and the acceptance
            angle in air.
      \item For a ray at that limiting angle, compute the geometric path
            length in $1.00\unit{m}$ of fibre and the time delay per kilometre
            relative to the axial ray. Why does this effect limit the bit rate
            of a step-index fibre, and what does a graded-index core do about
            it?
      \item The attenuation at $500\unit{nm}$ is $0.25\unit{dB/km}$. Compute
            the transmitted fraction after $10\unit{km}$. Then explain why
            silica fibres transmit infrared far better than violet, using the
            idea that absorption requires a photon energy matching an
            electronic transition.
      \end{enumerate}
\end{parts}

%======================================================================
\problem{10}{FY8 --- Quantisation, atoms, nuclei, radiation}{fy8}

\begin{parts}
\item A sodium lamp emits the doublet at $589.0\unit{nm}$ and
      $589.6\unit{nm}$.
      \begin{enumerate}[label=\roman*.]
      \item Compute the photon energies in $\unit{eV}$, their difference in
            $\unit{meV}$, and the frequency difference in $\unit{GHz}$.
      \item The light passes through a grating with $600$ lines per
            millimetre. Determine all diffraction maxima that exist and their
            angles, and how many lines must be illuminated to resolve the
            doublet in first order.
      \item Light of $500\unit{nm}$ falls on a metal with work function
            $2.28\unit{eV}$. Compute the threshold wavelength, the maximum
            kinetic energy of the photoelectrons and the stopping potential.
            Which of these quantities is independent of the light's intensity,
            and what does the intensity control instead?
      \end{enumerate}
\item The D--T fusion reaction is
      ${}^{2}\mathrm{H}+{}^{3}\mathrm{H}\to{}^{4}\mathrm{He}+n$, with masses
      $2.014102$, $3.016049$, $4.002602$ and $1.008665\unit{u}$, where
      $1\unit{u}=931.494\unit{MeV}/c^2$.
      \begin{enumerate}[label=\roman*.]
      \item Compute the energy released per reaction in MeV and, taking only
            the deuterium as fuel, the energy density in J/kg. Compare it
            with U-235 fission ($\approx200\unit{MeV}$ per nucleus) and with
            coal ($\approx30\unit{MJ/kg}$).
      \item A fusion plant produces $2.5\unit{GW}$ of fusion power. How many
            reactions per second does that require, and what mass of
            deuterium per year? Assuming seawater contains about
            $33\unit{g}$ of deuterium per cubic metre, estimate the volume of
            seawater needed per year, and judge whether fuel supply is the
            limiting engineering problem.
      \end{enumerate}
\item A ${}^{131}\mathrm{I}$ sample of activity $20.0\unit{GBq}$ is prepared
      for a patient; the half-life is $8.02\unit{d}$.
      \begin{enumerate}[label=\roman*.]
      \item What activity remains after $12.0\unit{h}$? Explain with
            $N=N_0e^{-\lambda t}$ why the number of nuclei in the sample does
            not visibly change over that time, and compute the initial number
            of nuclei.
      \item Why is ${}^{131}\mathrm{I}$ --- a $\beta^-$ and $\gamma$ emitter
            with an eight-day half-life --- well suited to thyroid imaging,
            whereas a pure $\alpha$ emitter of the same activity would be
            useless for it?
      \end{enumerate}
\item One patient absorbs $0.20\unit{Gy}$ from a gamma source distributed
      through the body; another absorbs the same $0.20\unit{Gy}$ from alpha
      emitters concentrated in lung tissue. Compute the equivalent dose in
      each case using $\mathrm{Sv}=\mathrm{Gy}\cdot w_R$, with $w_R=1$ for
      photons and $w_R=20$ for alpha particles, and judge which case carries
      the greater health risk in terms of \emph{equivalent} dose. Then state
      the further factor this conversion ignores and which makes the two cases
      differ even more.
\end{parts}

\end{document}
